\documentclass[10pt,a4paper]{amsart}
\usepackage[margin=19mm]{geometry}
\usepackage{amsmath,amssymb,mathtools}
\usepackage[scr=rsfs,cal=euler]{mathalfa}
\usepackage{xfrac}
\usepackage[unicode,hidelinks,bookmarks=false]{hyperref}
\newcommand{\ie}{i.e.\ }
\newcommand{\cf}{cf.\ }
\newcommand{\resp}{respectively\ }
\newcommand{\Subsection}{\subsection}
\providecommand{\nobreakdash}{\nobreak\hbox{-}}
\setlength{\emergencystretch}{2em}
\multlinegap0pt
\usepackage{bm}
\usepackage{bbm}
\usepackage{dsfont}
\usepackage{xspace}
\usepackage{cancel}
\usepackage{mathtools}
\usepackage[shortlabels]{enumitem}
\usepackage{amsthm}
\makeatletter
\@ifundefined{lemma}{\newtheorem{lemma}{Lemma}}{}
\@ifundefined{corollary}{\newtheorem{corollary}[lemma]{Corollary}}{}
\@ifundefined{proposition}{\newtheorem{proposition}[lemma]{Proposition}}{}
\@ifundefined{theorem}{\newtheorem{theorem}[lemma]{Theorem}}{}
\makeatother
\newcommand{\CE}{\operatorname{CE}}
\newcommand{\id}{\operatorname{id}}
\newcommand{\rC}{\mathrm{C}}
\newcommand{\ran}{\operatorname{ran}}
\title[$C^*$-supports and abnormalities of operator systems]{$C^*$-supports and abnormalities of operator systems}
\author{Rapha\"el Clou\^atre, Colin Krisko}
\date{Source: arXiv:2501.07544v2 (CC BY 4.0)}
\begin{document}
\maketitle

\noindent\emph{Source and licence.} $C^*$-supports and abnormalities of operator systems. Version arXiv:2501.07544v2 (CC BY 4.0), distributed under the Creative Commons Attribution 4.0 International licence, as recorded in the arXiv licence metadata for this exact version and shown on the abstract page https://arxiv.org/abs/2501.07544v2. Full rights notice: licenses/arxiv-2501.07544-CC-BY-4.0.txt.\par\smallskip

\noindent\textit{Editorial scope.} Counted unit: the Theorem whose source label is \texttt{T:Shtrivial} in the article (source0.tex lines 621--629, source label \texttt{T:Shtrivial}), delivered together with the article's own complete original proof of it (source0.tex lines 630--652, one \texttt{proof} environment of 23 lines). No mathematics is written, repaired or re-derived: statement and proof are the article's own text line for line. Required context retained uncounted: T:Hamana (lines 419--421); P:CEtrik (lines 468--474); P:CEcomp (lines 494--501); T:Cstarext (lines 562--568); C:coversenv (lines 593--599). Every definition, display and label of those lines is retained unchanged. Omitted from this package and not redistributed: 1--418, 422--467, 475--493, 502--561, 569--592, 600--620, 653--926. Nothing inside a retained excerpt is omitted or altered: every retained body is delivered exactly as the article's lines are written, so no omission receipt of any kind accompanies this package. Citations: the retained excerpts carry no citation command at all, so no bibliography entry of the article is transcribed and no reference is redistributed. Reference adaptation: none was needed and none was made; every reference inside the delivered unit resolves inside it, so no unresolved reference or citation is produced. Printed slips kept literal: no printed word, symbol, hypothesis or number of the counted statement or of its proof was altered. Layout and class: the article's own class is \texttt{amsart} with packages including graphicx, amscd, amsmath, amsthm, amsfonts, amssymb, mathrsfs, bm, enumerate, amsrefs, xcolor, hyperref, inputenc, tikz; the wrapper is the lane's pinned \texttt{amsart} editorial wrapper at 10pt on A4, which restates the theorem-like environments the retained text needs and adds the article's own macro definitions that the retained text needs (\texttt{\textbackslash CE}, \texttt{\textbackslash id}, \texttt{\textbackslash rC}, \texttt{\textbackslash ran}), with extra wrapper packages (bm, bbm, dsfont, xspace, cancel, mathtools, enumitem). The delivered numbering is the unit's own. Assets: no retained span contains a figure environment, a graphics-inclusion command, a PDF-page inclusion, a table or a TikZ picture; no image file is copied into this package, the package declares no asset dependency and no image file is redistributed with it. Licence: version arXiv:2501.07544v2 of this article is distributed under the Creative Commons Attribution 4.0 International licence, as recorded in the arXiv licence metadata for this exact version and shown on the abstract page https://arxiv.org/abs/2501.07544v2. A full rights notice is delivered at licenses/arxiv-2501.07544-CC-BY-4.0.txt. Characters: every retained excerpt is pure ASCII, so the delivered engine, XeLaTeX with its default Latin Modern text font, reports no missing character.
\begin{theorem}\label{T:Hamana}
Let $S\subset B(H)$ be an operator system and let $i:S\to B(H)$ denote the inclusion. Let $\theta:B(H)\to B(H)$ be an $S$-projection. Then, $\theta$ is minimal if and only if $( \ran \theta, i)$ is an injective envelope of $S$.   
\end{theorem}

\begin{proposition}\label{P:CEtrik}
Let $A$ and $B$ be unital $\rC^*$-algebras and let $\gamma:A\to B$ be a unital completely isometric map. Then, the following statements hold.
\begin{enumerate}[{\rm (i)}]
\item There is a unital $*$-homomorphism $\eta:\rC^*(\gamma(A))\to A$ such that $\eta\circ \gamma=\id$.
\item If  $\psi:\rC^*(\gamma(A))\to A$ is a unital completely positive map satisfing $\psi\circ \gamma=\id$ on $A$, then $\psi=\eta$. 
\end{enumerate} 
\end{proposition}

\begin{proposition}\label{P:CEcomp}
Let $S\subset B(H)$ be an operator system.  Let $\psi:B(H)\to B(H)$ be an $S$-projection and let $\theta:B(H)\to B(H)$ be a minimal $S$-projection. Then, the following statements hold.
\begin{enumerate}[{\rm(i)}]
\item The map \[\pi=\alpha_\theta\circ \theta\circ \alpha_\psi^{-1}|_{\CE(S,\psi)}: \CE(S,\psi)\to \CE(S,\theta)\] is a unital surjective $*$-homomorphism whose kernel is the Shilov ideal of $\alpha_\psi(S)$ in $\CE(S,\psi)$.
\item  If $\alpha_\psi^{-1}(\CE(S,\psi))$ is contained in the range of $\theta$, then \[\alpha_\psi^{-1}(\CE(S,\psi))=\alpha_\theta^{-1}( \CE(S,\theta)).\]
\item Assume that there is a unital $*$-homomorphism $\sigma:\CE(S,\theta)\to \CE(S,\psi)$ such that $\sigma\circ\alpha_\theta|_S=\alpha_\psi|_S$. Then, there is a minimal $S$-projection $\theta':B(H)\to B(H)$ such that $\alpha_{\theta'}^{-1}(\CE(S,\theta'))=\alpha_\psi^{-1}(\CE(S,\psi))$. 
\end{enumerate}
\end{proposition}

\begin{theorem}\label{T:Cstarext}
Let $S\subset B(H)$  be an operator system. Let $X\subset B(H)$ be an operator system containing $S$. Then, the following statements are equivalent.
\begin{enumerate}[{\rm (i)}]
\item $X$ is a $\rC^*$-support of $S$.
\item There is an $S$-projection $\theta:B(H)\to B(H)$ such that  $X=\alpha_\theta^{-1}(\CE(S,\theta))$.
\end{enumerate}
\end{theorem} 

\begin{corollary}\label{C:coversenv}
Let $S\subset B(H)$  be an operator system. Let $X\subset B(H)$ be an operator system containing $S$.  Then, the following statements are equivalent.
\begin{enumerate}[{\rm (i)}]
\item If $(A,j)$ is a $\rC^*$-envelope of $S$, then there is a unital, surjective, completely isometric map  $\Gamma:A\to X$ such $\Gamma\circ j=\id$. 
\item There is a minimal $S$-projection $\theta:B(H)\to B(H)$ such that  \[X=\alpha_\theta ^{-1}(\CE(S,\theta)).\]
\end{enumerate}
\end{corollary} 

\begin{theorem}\label{T:Shtrivial}
Let $S\subset B(H)$  be an operator system. Then, the following statements are equivalent.
\begin{enumerate}[{\rm (i)}]
\item If $X\subset B(H)$ is a $\rC^*$-support of $S$ and $(A,j)$ is a $\rC^*$-envelope of $S$, then there is a unital, surjective, completely isometric map $\Gamma:A\to X$ such that $\Gamma\circ j=\id$. 
\item $(\rC^*(S),i)$ is a  $\rC^*$-envelope of $S$, where $i:S\to B(H)$ is the inclusion map.
\item $\rC^*(S)$ is contained in some injective operator system $R\subset B(H)$  such that $(R,\kappa)$ is an injective envelope of $S$, where $\kappa:R\to B(H)$ is the inclusion map.
\item For each $a\in \rC^*(S)$, there is an injective operator system $R\subset B(H)$  such that $a\in R$ and $(R,\kappa)$ is an injective envelope of $S$, where $\kappa:R\to B(H)$ is the inclusion map.
\end{enumerate}
\end{theorem}

\begin{proof}
(i) $\Rightarrow$ (ii): Let $(A,j)$ be a $\rC^*$-envelope of $S$. Plainly,  $(\rC^*(S),i)$ is a $\rC^*$-extension of $S$, and hence a $\rC^*$-support. By assumption, there is a unital, surjective, completely isometric map $\Gamma:A\to \rC^*(S)$ such that $\Gamma\circ j=\id$. By Proposition \ref{P:CEtrik}, we infer that $\Gamma$ is a $*$-isomorphism. Thus, $(\rC^*(S),i)$ is a  $\rC^*$-envelope of $S$.


(ii) $\Rightarrow$ (i):  Fix a $\rC^*$-support $X\subset B(H)$ of $S$. By Theorem \ref{T:Cstarext}, there is an $S$-projection $\psi:B(H)\to B(H)$ such that $X=\alpha_\psi^{-1}(\CE(S,\psi))$. Furthermore, denoting the range of $\psi$ by $R_\psi\subset B(H)$, we know that $\alpha_\psi\circ \psi:\rC^*(R_\psi)\to \CE(R_\psi,\psi)$ is a unital  $*$-homomorphism by construction of the Choi--Effros product. Let $\pi=(\alpha_\psi\circ \psi)|_{\rC^*(S)}$, which is a $*$-homomorphism satisfying $\pi|_S=\alpha_\psi|_S$. In particular, the image of $\pi$ is $\CE(S,\psi)$. 

Next, our assumption is that $(\rC^*(S),i)$ is a $\rC^*$-envelope of $S$, so there must exist a unital surjective $*$-homomorphism $\sigma:\CE(S,\psi)\to \rC^*(S)$ such that $\sigma \circ \alpha_\psi|_S=\id$. We have
\[
\sigma\circ \pi|_S=\sigma\circ \alpha_\psi|_S=\id 
\]
and
\[
\pi\circ \sigma\circ \alpha_\psi|_S=\pi|_S=\alpha_\psi|_S.
\]
Consequently, because $\pi$ and $\sigma$ are $*$-homomorphisms, we find $\pi=\sigma^{-1}$ so $\pi$ is completely isometric.
Hence, we may define $\Gamma=\alpha_\psi^{-1}\circ \pi:\rC^*(S)\to X$, which is a unital, surjective, completely isometric map satisfying $\Gamma|_S=\id$.  This establishes the desired statement for the $\rC^*$-envelope $(\rC^*(S),i)$. The corresponding statement for any $\rC^*$-envelope of $S$ follows readily by uniqueness.

(ii) $\Rightarrow$ (iii): We know that $(\rC^*(S),i)$ is a $\rC^*$-envelope of $S$. Thus, applying Corollary \ref{C:coversenv}, we find a minimal $S$-projection $\theta:B(H)\to B(H)$ such that  \[\rC^*(S)=\alpha_\theta ^{-1}(\CE(S,\theta)).\] In particular, this shows that $\rC^*(S)$ is contained in the range of $\theta$. Furthermore, if we let $\kappa:\ran\theta\to B(H)$ denote the inclusion, then $(\ran \theta, \kappa)$  is an injective envelope of $S$ by Theorem \ref{T:Hamana}.

(iii) $\Rightarrow$ (iv): This is trivial. 

(iv) $\Rightarrow$ (ii): Let $a\in \rC^*(S)$ be an element in the Shilov ideal of $S$ in $\rC^*(S)$. By assumption, there is an injective operator system $R\subset B(H)$  containing $a$ such that $(R,\kappa)$ is an injective envelope of $S$, where $\kappa:R\to B(H)$ denote the inclusion map. Since $R$ is injective, we may choose an $S$-projection $\theta:B(H)\to B(H)$ with range $R$. In turn, Theorem \ref{T:Hamana} implies that $\theta$ must be a minimal $S$-projection.  By virtue of Proposition \ref{P:CEcomp}(i) applied with $\psi=\id$, we see that $a$ must lie in the kernel of $\alpha_\theta\circ \theta$, and hence $\theta(a)=0$. On other hand, $\theta(a)=a$ since $a$ lies in the range of the idempotent $\theta$, so we conclude $a=0$. Therefore, the Shilov ideal of $S$ in $\rC^*(S)$ is trivial, so that $(\rC^*(S),i)$ is a  $\rC^*$-envelope of $S$.
\end{proof}

\end{document}
